\documentclass[11pt]{article}
\newcommand{\SheetCode}{Theory 07}
\usepackage[margin=25mm]{geometry}
\usepackage[T1]{fontenc}
\usepackage{lmodern}
\DeclareUnicodeCharacter{2605}{$\star$}
\usepackage{microtype}
\usepackage{amsmath,amssymb,mathtools,bm}
\usepackage{enumitem}
\usepackage{xcolor}
\usepackage{fancyhdr}
\usepackage{hyperref}
\usepackage[most]{tcolorbox}
\definecolor{agiBlue}{HTML}{173B57}
\definecolor{agiGold}{HTML}{B7791F}
\definecolor{agiLight}{HTML}{EFF6FA}
\definecolor{agiGreen}{HTML}{184D3B}
\hypersetup{colorlinks=true,linkcolor=agiBlue,urlcolor=agiBlue}
\setlength{\parindent}{0pt}
\setlength{\parskip}{0.55em}
\setlength{\headheight}{14pt}
\setlist{nosep,leftmargin=*}
\pagestyle{fancy}
\fancyhf{}
\fancyhead[L]{\small Part of the \href{https://github.com/mmikhasenko/GIModel.jl}{Agentic Gottried--Isgur} project}
\fancyhead[R]{\SheetCode}
\fancyfoot[C]{\thepage}
\newcommand{\dd}{\,\mathrm d}
\newcommand{\ket}[1]{\lvert #1\rangle}
\newcommand{\bra}[1]{\langle #1\rvert}
\newcommand{\braket}[2]{\langle #1\mid #2\rangle}
\newcommand{\expect}[1]{\left\langle #1\right\rangle}
\newcommand{\op}[1]{\widehat{#1}}
\newcommand{\GeV}{\,\mathrm{GeV}}
\newcommand{\R}{\mathbb{R}}
\newtcolorbox{goals}{colback=agiLight,colframe=agiBlue,title=Learning outcomes}
\newtcolorbox{reading}{colback=white,colframe=agiGold,title=Book reference,
  after skip=5mm}
\newtcolorbox{paperbridge}{colback=agiGreen!7,colframe=agiGreen,
  colbacktitle=agiGreen,coltitle=white,title=Connection to the GI paper}
\newtcolorbox{problem}[1]{colback=white,colframe=agiBlue,title={#1},
  before skip=3.5mm,after skip=1mm}
\newtcolorbox{solution}{colback=black!2,colframe=black!45,title=Solution}
\newtcolorbox{quiz}{colback=white,colframe=agiGold,title=Post-solution concept check}
\newtcolorbox{checkpoint}{colback=white,colframe=agiGold,title=Interpretation checkpoint}
\newcommand{\sheettitle}[3]{%
  \begin{center}
  {\Large\bfseries #1}\par
  \vspace{0.2em}{\color{agiBlue}#2}\par
  \vspace{0.4em}\small #3
  \end{center}\vspace{0.5em}}
\newcommand{\difficulty}[1]{\textbf{Difficulty:} #1}
\newcommand{\timebox}[1]{\textbf{Suggested time:} #1}
\newcommand{\answer}{\par\smallskip\color{black!50}\textbf{Answer:} }

\begin{document}
\sheettitle{Sheet 7: When spectroscopic labels stop being states}{Two-level mixing, open flavor, and isoscalar annihilation}{Prerequisite: Sheets 1--6. \difficulty{★★/★★★} \quad \timebox{2--3 hours}}

\begin{goals}
Diagonalize and interpret a $2\times2$ mass matrix; track phase-invariant
observables; identify antisymmetric spin--orbit and tensor mixing; build an
isoscalar flavor block; understand why mixed states need component waves.
\end{goals}
\begin{reading}
Selected reading in Landau--Lifshitz, \emph{Quantum Mechanics: Non-Relativistic
Theory}: Ch. VI, ``Perturbation Theory,'' Ch. XII, ``The Theory of Symmetry,''
and Ch. XIV, ``Addition of Angular Momenta.'' These are selected references;
the sheet develops the mixing problem in its own order.
\end{reading}


Here \emph{open flavor} means that the quark and antiquark have different
flavors, so charge conjugation does not return the same meson. A \emph{basis
phase} is the arbitrary choice of multiplying any basis ket by $-1$ (or, more
generally, by a complex phase); observables cannot depend on that choice.

\begin{problem}{Problem 1 ★ — Solve the universal mixing matrix}
In an orthonormal basis $\ket a,\ket b$, the mass matrix is
$M=\begin{pmatrix}M_a&V\\V&M_b\end{pmatrix}$, with real energies
$M_a,M_b,V$. Find its ordered eigenvalues $M_-\le M_+$ and normalized
states in the convention
$\ket-=\cos\theta\ket a-\sin\theta\ket b$,
$\ket+=\sin\theta\ket a+\cos\theta\ket b$.
Specify the angle branch, including $V=0$ and $M_a=M_b$.

\textbf{Hint:} Solve $\det(M-\lambda I)=0$, then require
$\bra- M\ket+=0$ and check which diagonal entry is lower.
\end{problem}
\begin{solution}
\[
M_\pm=\frac{M_a+M_b}{2}\pm
\sqrt{\left(\frac{M_a-M_b}{2}\right)^2+V^2},\qquad
\tan2\theta=\frac{2V}{M_b-M_a}.
\]
One branch that assigns $\ket-$ to the lower eigenvalue is
$2\theta=\operatorname{atan2}(2V,M_b-M_a)$, where $\operatorname{atan2}(y,x)$
is the quadrant-aware angle of $(x,y)$. For $V=0$, take $\theta=0$ if
$M_a<M_b$ and $\theta=\pi/2$ if $M_a>M_b$. If $M_a=M_b$ and $V\ne0$,
take $\theta=\operatorname{sgn}(V)\pi/4$. When both $V=0$ and $M_a=M_b$,
any orthonormal pair is an eigenbasis.
The sign of $\theta$ depends on basis phases; eigenvalues and squared component
magnitudes do not.
\end{solution}
\begin{quiz}
If $\ket b\to-\ket b$, what measurable quantity changes? \answer none; $V$ and
$\theta$ change sign together, while masses and probabilities remain fixed.
\end{quiz}

\begin{problem}{Problem 2 ★★ — Prove level repulsion}
Use the eigenvalues of Problem 1 to prove that $V\ne0$ implies
$M_-<\min(M_a,M_b)$ and $M_+>\max(M_a,M_b)$.
For $M_a<M_b$, define $\Delta=M_b-M_a>0$ and obtain
$M_--M_a$ and $M_+-M_b$ through order $V^2/\Delta$
when $|V|\ll\Delta$.

\textbf{Hint:} Factor $\Delta/2$ out of the square root and expand in
$4V^2/\Delta^2$.
\end{problem}
\begin{solution}
The square root exceeds $|M_a-M_b|/2$, proving both inequalities. If
$M_a<M_b$, expansion gives
$M_-=M_a-V^2/(M_b-M_a)+O(V^4)$ and
$M_+=M_b+V^2/(M_b-M_a)+O(V^4)$.
\end{solution}
\begin{quiz}
Why can a small off-diagonal interaction cause large mixing? \answer the angle
depends on the ratio $2V/(M_b-M_a)$; near degeneracy the denominator is small.
\end{quiz}

\begin{problem}{Problem 3 ★★ — Locate antisymmetric spin--orbit mixing}
For positive constituent masses $m_1,m_2$, express
$\bm S_1\cdot\bm L/m_1^2+\bm S_2\cdot\bm L/m_2^2$ as
$C_+(\bm S_1+\bm S_2)\cdot\bm L+C_-(\bm S_1-\bm S_2)\cdot\bm L$.
Find $C_+,C_-$ and evaluate them at $m_1=m_2$.
Identify which spin operator can connect a singlet to a triplet.

\textbf{Hint:} Match the coefficients of $\bm S_1\cdot\bm L$ and
$\bm S_2\cdot\bm L$. Singlets and triplets have opposite spin-exchange symmetry.
\end{problem}
\begin{solution}
\[
\frac12\left(\frac1{m_1^2}+\frac1{m_2^2}\right)
(\bm S_1+\bm S_2)\cdot\bm L+
\frac12\left(\frac1{m_1^2}-\frac1{m_2^2}\right)
(\bm S_1-\bm S_2)\cdot\bm L.
\]
The antisymmetric coefficient vanishes at $m_1=m_2$. Its spin operator connects
$^1L_L$ and $^3L_L$ in unequal-mass mesons.
\end{solution}
\begin{quiz}
Why is open flavor the natural laboratory for this mixing? \answer unequal
constituent masses activate the antisymmetric coefficient and charge
conjugation is not a prohibiting quantum number.
\end{quiz}

\begin{problem}{Problem 4 ★★ — Build the isoscalar flavor basis}
The states $\ket{u\bar u},\ket{d\bar d},\ket{s\bar s}$ are orthonormal.
Assume $\bra{f\bar f}H_A\ket{g\bar g}=A$ for every
$f,g\in\{u,d,s\}$, where $A$ is a real energy.
With $\ket n=(\ket{u\bar u}+\ket{d\bar d})/\sqrt2$ and
$\ket s=\ket{s\bar s}$, compute all four entries of $H_A$ in the $(n,s)$
basis. Mass and radial-wave factors are omitted.

\textbf{Hint:} Expand each $n$ bra and ket before summing amplitudes.
\end{problem}
\begin{solution}
Every elementary flavor transition has amplitude $A$. Coherent sums give
\[
H_A=A\begin{pmatrix}2&\sqrt2\\\sqrt2&1\end{pmatrix}.
\]
Its eigenvalues are $0$ and $3A$: one flavor combination cancels and the singlet
combination adds coherently. Real GI elements include channel-dependent masses,
$\alpha_s$, and origin-smeared waves.
\end{solution}
\begin{quiz}
Why can annihilation shift one isoscalar strongly while leaving an orthogonal
combination nearly untouched? \answer amplitudes interfere coherently in flavor
space.
\end{quiz}

\begin{problem}{Problem 5 ★★★ — Check the GI annihilation dimensions}
For
$A_{ji}\propto A_0[\alpha_s(M_j^2)\alpha_s(M_i^2)]^{n/2}
S_L(\Psi_j)S_L(\Psi_i)/(m_im_j)$, suppose
$[S_L]=E^{L+3/2}$. Here $S_L(\Psi)$ is a linear operation that turns the whole
wavefunction into one number by weighting its short-distance momentum content.
The coupling $\alpha_s$ and exponent $n$ are dimensionless;
$M_i,m_i$ have energy dimension, and the omitted proportionality factor
is dimensionless. The two states have the same $L$.
Determine the dimension required for $A_0$ so that
$A_{ji}$ is an energy. Explain why the near-origin wave is decisive.
\end{problem}
\begin{solution}
The two $S_L$ factors contribute $E^{2L+3}$ and the masses contribute $E^{-2}$,
so $A_0$ must have dimension $E^{-2L}$ to leave $E^1$. The origin-smeared
functional selects short-distance and derivative information: two states with
similar masses but different nodes or $L$ can have very different annihilation
strengths.
\end{solution}
\begin{quiz}
Why is it especially risky here to represent the wave only at widely separated
radii? \answer eigenvalues depend on the wave over the whole radial range, but
the annihilation quantity gives exceptional weight to the small region near the
origin and therefore magnifies local errors there.
\end{quiz}

\begin{problem}{Problem 6 ★★★ — Reconstruct a mixed observable}
Let $\ket\Psi=c_a\ket{a,\psi_a}+c_b\ket{b,\psi_b}$ and let $O$ be diagonal in
the discrete labels. Derive $\expect O$. Why is assigning one ``representative
wave'' to $\ket\Psi$ generally wrong?

\textbf{Hint:} Take orthogonal discrete states $\ket a,\ket b$,
normalized radial waves, and $|c_a|^2+|c_b|^2=1$.
Write $O=\ket a\bra a\otimes O_a+\ket b\bra b\otimes O_b$.
The radial functions need not be equal or mutually orthogonal.
\end{problem}
\begin{solution}
If $O$ has no cross-label element,
$\expect O=|c_a|^2\bra{\psi_a}O_a\ket{\psi_a}
+|c_b|^2\bra{\psi_b}O_b\ket{\psi_b}$. With allowed cross terms, coherent
$c_a^*c_b$ contributions must also be retained. A single wave discards the
label-dependent radial functions and cannot reproduce both cases. The physical
object is the list of coefficients paired with each component's radial wave in
the representation in which that wave was calculated.
\end{solution}
\begin{quiz}
What data must survive a mixing stage? \answer eigenvalues, normalized mixing
coefficients, basis-state identities, and access to each component's radial
wave. Keeping only a mass and one chosen wave would discard information needed
for later observables.
\end{quiz}

\begin{paperbridge}
The paragraph following GI Eq. (14) defines the second-stage tensor
$^3L_J\leftrightarrow{}^3L'_J$ and unequal-mass antisymmetric spin--orbit
$^3L_J\leftrightarrow{}^1L_J$ mass matrices. Eqs. (15)--(17) define the
isoscalar annihilation matrix and its short-distance wavefunction factor;
Eqs. (18a)--(18b) give the special pseudoscalar alternatives. Physical
isoscalar compositions are tabulated in Table III.
\end{paperbridge}
\end{document}
